\documentclass{report}
\usepackage{amsfonts, amsmath}
\newcommand\R{{\mathbb R}}
\newcommand\Q{{\mathbb Q}}
\newcommand\N{{\mathbb N}}
\pagestyle{empty}
\begin{document}
\large
\centerline{\Huge \bf Analisi Matematica I}
\vskip.2cm
\centerline{\Huge Secondo Esonero}
\renewcommand{\thefootnote}{ } 
\footnote{\sl Avete 3 ore di tempo. Ogni esercizio
vale sei punti. La sufficienza si ottiene con un punteggio $\geq$
18. {\bf Solo le risposte chiaramente giustificate saranno prese in
considerazione.} {\scshape Elaborati illeggibili o confusi verranno ignorati.}
{\sffamily La
sintesi sar\`a particolarmente apprezzata.}}
\renewcommand{\thefootnote}{\arabic{footnote}} 

\vskip1cm
\begin{enumerate}
\item Si consideri la funzione
\[
f(x)=x^4-\sin x.
\]
Se ne faccia il grafico e si dica quante sono le soluzioni
dell'equazione $f(x)=0$.

\item Tra tutti i triangoli di vertici $(0,0)$, $(0,1)$ $(x,1-x^2)$,
per $x\in[-1,1]$, si trovi quello di area massima.

\item Si dimostri che
\[
x^2\geq 2(1-\cos x).
\]
\item  Si disegni l'insieme $A=\{(x,y)\in\R^2\;|\; x^2-y^2>0\}$.
\item Si dia un esempio di funzione definita su tutto $\R$ che sia
continua ma non derivabile in due punti.
\item Si determini, senza usare il calcolatore, il numero delle
soluzioni dell'equazione
\[
2e^x-3=0
\]
nell'intervallo $x\in (0,.5)$.

\end{enumerate}
\end{document}
